\documentclass[12pt,a4paper]{article}
\usepackage{geometry}
\geometry{left=2.5cm,right=2.5cm,top=2.0cm,bottom=2.5cm}
\usepackage[english]{babel}
\usepackage{amsmath,amsthm}
\usepackage{amsfonts}
\usepackage[longend,ruled,linesnumbered]{algorithm2e}
\usepackage{fancyhdr}
\usepackage{ctex}
\usepackage{array}
\usepackage{listings}
\usepackage{color}
\usepackage{graphicx}

\begin{document}


\title{
{《优化理论与算法》第 5 次作业
}
}
\date{}

\author{
姓名：\underline{焦传斌}~~~~~~
学号：\underline{2024111223}~~~~~~}

\maketitle

\begin{itemize}
	\item 作业截止于{\bf 周五（4/10/2026）}，在Canvas系统中提交PDF或照片文件
	\item 作业题目的tex文件可以在Canvas中下载
	\item 鼓励相互讨论，但不要抄袭.
\end{itemize}
%%%%%%%%%%%%%%%%%%%%%%%%%%%%%%%%%%%%%%%%%%%%%%%%%%%%%%%%%%%%%%%%%
%%%%%%%%%%%%%%%%%%%%%%%%%%%%%%%%%%%%%%%%%%%%%%%%%%%%%%%%%%%%%%%%%

\section{阅读与积累}
\paragraph{1.1}
\begin{itemize}
	\item 阅读Bertsimas and Tsitsikils，也就是Introduction to Linear Optimization教材，第4.1-4.4节、5.1-5.2节相关内容。
\end{itemize}

\section{习题}

\paragraph{2.1}
考虑以下线性规划问题：
\begin{align*}
\min \;&\quad x_1 - x_2 \\
\text{s.t.} \;&\quad2x_1 + 3x_2 - x_3 + x_4 \leq 0 \\
 \;&\quad3x_1 + x_2 + 4x_3 - 2x_4 \geq 3 \\
 \;&\quad-x_1 - x_2 + 2x_3 + x_4 = 6 \\
 \;&\quad x_1 \leq 0, x_2 \text{ 无约束 }, x_3 \geq 0, x_4 \text{ 无约束 }
\end{align*}

写出相应的对偶问题。

解：

\begin{align*}
\max\;&\quad 3y_2+6y_3\\
\text{s.t. }&\quad 2y_1+3y_2-y_3\ge 1\\
&\quad 3y_1+y_2-y_3=-1\\
&\quad -y_1+4y_2+2y_3\le 0\\
&\quad y_1-2y_2+y_3=0\\
&\quad y_1\le 0,\ y_2\ge 0,\ y_3\text{ 无约束}.
\end{align*}

\paragraph{2.2}
考虑以下线性规划问题：

\begin{align*}
   \max \quad &\; 3x_1 + 2x_2\\
   s.t.\quad &\;x_1 + x_2 \leq 4 \\
&\;2x_1 + x_2 \leq 5 \\
&\;x_1, x_2 \geq 0
\end{align*}
利用互补松弛条件判断下列解是否是最优解。

解：其对偶问题为
\begin{align*}
\min\;&\quad 4y_1+5y_2\\
\text{s.t. }&\quad y_1+2y_2\ge 3\\
&\quad y_1+y_2\ge 2\\
&\quad y_1,y_2\ge 0.
\end{align*}
互补松弛条件为
\begin{align*}
y_1(4-x_1-x_2)&=0,\quad y_2(5-2x_1-x_2)=0,\\
x_1[(y_1+2y_2)-3]&=0,\quad x_2[(y_1+y_2)-2]=0.
\end{align*}

\begin{itemize}
    \item 对于解 \(x_1=1,\ x_2=3\)：该点原问题可行，由互补松弛条件有
    \[
    y_1+2y_2=3,\qquad y_1+y_2=2.
    \]
    解得
    \[
    y_1=1,\qquad y_2=1.
    \]
    该解满足对偶可行性条件，故存在与之对应的对偶可行解满足互补松弛条件，因此 \((1,3)\) 是最优解。

    \item 对于解 \(x_1=2,\ x_2=1\)：该点原问题可行，
由互补松弛条件知 
    \[
	y_1=0，\qquad
    y_1+2y_2=3,\qquad y_1+y_2=2.
    \]
无解，所以 \((2,1)\) 不是最优解。
\end{itemize}

\paragraph{2.3}
试证明：线性系统$Ax=b, x \geq 0$无解，只需找到某个向量$y$使之满足$A^Ty\geq0, b^Ty<0$。

证明：若存在向量 $y$ 使得
\[
A^Ty\ge 0,\qquad b^Ty<0,
\]

设存在 $x\ge 0$ 使得 $Ax=b$，则有
\[
b^Ty=(Ax)^Ty=x^TA^Ty.
\]
由于 $x\ge 0$ 且 $A^Ty\ge 0$，故
\[
x^TA^Ty\ge 0,
\]
从而
\[
b^Ty\ge 0.
\]
这与已知条件 $b^Ty<0$ 矛盾。

因此，线性系统 $Ax=b,\ x\ge 0$ 无解。


\paragraph{2.4}

先化标准型：

\[
\min w=-2x_1-4x_2-x_3-x_4
\]

s.t.
\[
\begin{aligned}
x_1+3x_2+x_4+x_5&=4,\\
2x_1+x_2+x_6&=3,\\
x_2+4x_3+x_4+x_7&=3,
\end{aligned}
\qquad x_i\ge 0\ (i=1,\dots,7)
\]






\textbf{(1)}

初始表：

\[
\begin{array}{c|rrrrrrr|r}
B & x_1&x_2&x_3&x_4&x_5&x_6&x_7&b\\
\hline
  & -2&-4&-1&-1&0&0&0&0\\
x_5&1&3&0&1&1&0&0&4\\
x_6&2&1&0&0&0&1&0&3\\
x_7&0&1&4&1&0&0&1&3
\end{array}
\]

 $x_2$ 入基。
比值检验：$\frac{4}{3},\ 3,\ 3$，故 $x_5$ 出基。

第一次迭代后：

\[
\begin{array}{c|rrrrrrr|r}
B & x_1&x_2&x_3&x_4&x_5&x_6&x_7&b\\
\hline
  & -\frac23&0&-1&\frac13&\frac43&0&0&\frac{16}{3}\\
x_2&\frac13&1&0&\frac13&\frac13&0&0&\frac43\\
x_6&\frac53&0&0&-\frac13&-\frac13&1&0&\frac53\\
x_7&-\frac13&0&4&\frac23&-\frac13&0&1&\frac53
\end{array}
\]

$x_3$ 入基。
比值检验：只有第三行该列为正，故 $x_7$ 出基。

第二次迭代后：

\[
\begin{array}{c|rrrrrrr|r}
B & x_1&x_2&x_3&x_4&x_5&x_6&x_7&b\\
\hline
  & -\frac34&0&0&\frac12&\frac54&0&\frac14&\frac{23}{4}\\
x_2&\frac13&1&0&\frac13&\frac13&0&0&\frac43\\
x_6&\frac53&0&0&-\frac13&-\frac13&1&0&\frac53\\
x_3&-\frac1{12}&0&1&\frac16&-\frac1{12}&0&\frac14&\frac5{12}
\end{array}
\]

 $x_1$ 入基。
比值检验：
\[
\frac{4/3}{1/3}=4,\qquad \frac{5/3}{5/3}=1
\]
故 $x_6$ 出基。

最终表：

\[
\begin{array}{c|rrrrrrr|r}
B & x_1&x_2&x_3&x_4&x_5&x_6&x_7&b\\
\hline
  & 0&0&0&\frac7{20}&\frac{11}{10}&\frac9{20}&\frac14&\frac{13}{2}\\
x_2&0&1&0&\frac25&\frac25&-\frac15&0&1\\
x_1&1&0&0&-\frac15&-\frac15&\frac35&0&1\\
x_3&0&0&1&\frac3{20}&-\frac1{10}&\frac1{20}&\frac14&\frac12
\end{array}
\]

因为所有检验数都满足 $\bar c_j\ge 0$，故达到最优。

令非基变量
\[
x_4=x_5=x_6=x_7=0
\]
得
\[
x_1=1,\quad x_2=1,\quad x_3=\frac12,\quad x_4=0
\]
故
\[
w^*=-\frac{13}{2}
\]
原最大化问题最优值为
\[
z^*=\frac{13}{2}
\]

结论：
\[
(x_1,x_2,x_3,x_4)=\left(1,1,\frac12,0\right),\qquad z^*=\frac{13}{2}
\]

\textbf{（2）在最优基下讨论 $b=[4,3,3]^T$ 的变化范围}

最终基为
\[
\{x_2,x_1,x_3\}
\]
因此最终表中 $x_5,x_6,x_7$ 三列，就是 $B^{-1}$ 的三列。
当前基变量值为
\[
(x_2,x_1,x_3)^T=(1,1,\tfrac12)^T
\]

1）$b_1$ 单独变化

设 $b_1=4+\lambda$，则
\[
\begin{pmatrix}x_2\\x_1\\x_3\end{pmatrix}
=
\begin{pmatrix}1\\1\\\frac12\end{pmatrix}
+\lambda
\begin{pmatrix}\frac25\\-\frac15\\-\frac1{10}\end{pmatrix}
\]
故
\[
1+\frac25\lambda\ge 0,\quad 1-\frac15\lambda\ge 0,\quad \frac12-\frac1{10}\lambda\ge 0
\]
解得
\[
-\frac52\le \lambda\le 5
\]
所以
\[
b_1\in \left[\frac32,9\right]
\]

2）$b_2$ 单独变化

设 $b_2=3+\lambda$，则
\[
\begin{pmatrix}x_2\\x_1\\x_3\end{pmatrix}
=
\begin{pmatrix}1\\1\\\frac12\end{pmatrix}
+\lambda
\begin{pmatrix}-\frac15\\\frac35\\\frac1{20}\end{pmatrix}
\]
故
\[
1-\frac15\lambda\ge 0,\quad 1+\frac35\lambda\ge 0,\quad \frac12+\frac1{20}\lambda\ge 0
\]
解得
\[
-\frac53\le \lambda\le 5
\]
所以
\[
b_2\in \left[\frac43,8\right]
\]

3）$b_3$ 单独变化

设 $b_3=3+\lambda$，则
\[
\begin{pmatrix}x_2\\x_1\\x_3\end{pmatrix}
=
\begin{pmatrix}1\\1\\\frac12\end{pmatrix}
+\lambda
\begin{pmatrix}0\\0\\\frac14\end{pmatrix}
\]
故只需
\[
\frac12+\frac14\lambda\ge 0
\]
即
\[
\lambda\ge -2
\]
所以
\[
b_3\in [1,+\infty)
\]

结论：
\[
b_1\in \left[\frac32,9\right],\quad b_2\in \left[\frac43,8\right],\quad b_3\in [1,+\infty)
\]

\textbf{（3）在最优基下讨论原目标系数 $c=[2,4,1,1]^T$ 的变化范围}

因为原问题是
\[
\max z=2x_1+4x_2+x_3+x_4
\]
而标准型里做的是
\[
\min w=-2x_1-4x_2-x_3-x_4
\]
所以若原问题某个系数增加 $\lambda$，则 minimize 标准型里对应系数减少 $\lambda$。

当前非基变量是
\[
x_4,x_5,x_6,x_7
\]
要保持当前基解仍最优，只需保持它们的新检验数仍满足 $\bar c_j\ge 0$。

1）$c_1$ 变化

设原问题 $c_1=2+\lambda$，则标准型中第1个系数变为 $-2-\lambda$。
由最终表中 $x_1$ 行可得：
\[
\bar c_4=\frac7{20}-\frac15\lambda\ge 0
\]
\[
\bar c_5=\frac{11}{10}-\frac15\lambda\ge 0
\]
\[
\bar c_6=\frac9{20}+\frac35\lambda\ge 0
\]
解得
\[
-\frac34\le \lambda\le \frac74
\]
所以
\[
c_1\in \left[\frac54,\frac{15}4\right]
\]

2）$c_2$ 变化

设原问题 $c_2=4+\lambda$。
由最终表中 $x_2$ 行可得：
\[
\bar c_4=\frac7{20}+\frac25\lambda\ge 0
\]
\[
\bar c_5=\frac{11}{10}+\frac25\lambda\ge 0
\]
\[
\bar c_6=\frac9{20}-\frac15\lambda\ge 0
\]
解得
\[
-\frac78\le \lambda\le \frac94
\]
所以
\[
c_2\in \left[\frac{25}8,\frac{25}4\right]
\]

3）$c_3$ 变化

设原问题 $c_3=1+\lambda$。
由最终表中 $x_3$ 行可得：
\[
\bar c_4=\frac7{20}+\frac3{20}\lambda\ge 0
\]
\[
\bar c_5=\frac{11}{10}-\frac1{10}\lambda\ge 0
\]
\[
\bar c_6=\frac9{20}+\frac1{20}\lambda\ge 0
\]
\[
\bar c_7=\frac14+\frac14\lambda\ge 0
\]
解得
\[
-1\le \lambda\le 11
\]
所以
\[
c_3\in [0,12]
\]

4）$c_4$ 变化

设原问题 $c_4=1+\lambda$，则标准型中非基变量 $x_4$ 的系数变为 $-1-\lambda$。

故其检验数变为
\[
\bar c_4=\frac7{20}-\lambda\ge 0
\]
解得
\[
\lambda\le \frac7{20}
\]
所以
\[
c_4\in \left(-\infty,\frac{27}{20}\right]
\]

结论：
\[
c_1\in \left[\frac54,\frac{15}4\right],\quad c_2\in \left[\frac{25}8,\frac{25}4\right],\quad c_3\in [0,12],\quad c_4\in \left(-\infty,\frac{27}{20}\right]
\]

\paragraph{2.5}


原问题
化为
\[
\min w=-x_1-2x_2
\]

s.t.
\[
\begin{aligned}
x_1+s_1&=100,\\
2x_2+s_2&=200,\\
x_1+x_2+s_3&=150,
\end{aligned}
\qquad x_1,x_2,s_1,s_2,s_3\ge 0
\]

原问题的最优表可写为
\[
\begin{array}{c|rrrrr|r}
B & x_1&x_2&s_1&s_2&s_3&b\\
\hline
  & 0&0&0&\frac12&1&250\\
s_1&0&0&1&\frac12&-1&50\\
x_2&0&1&0&\frac12&0&100\\
x_1&1&0&0&-\frac12&1&50
\end{array}
\]
即原最优解为
\[
x_1=50,\ x_2=100,\ w^*=-250
\]
也就是原最大化问题最优值
\[
z^*=250
\]

\textbf{（1）}

新问题：
\[
\max z=x_1+2x_2+3x_3
\]
化为标准型：
\[
\min w=-x_1-2x_2-3x_3
\]

s.t.
\[
\begin{aligned}
x_1+x_3+s_1&=100,\\
2x_2+x_3+s_2&=200,\\
x_1+x_2+x_3+s_3&=150,
\end{aligned}
\qquad x_1,x_2,x_3,s_1,s_2,s_3\ge 0
\]





在原最优表基础上加入新列。
因为新列 $a_3=(1,1,1)^T$，在当前基下
\[
B^{-1}a_3=\left(\frac12,\frac12,\frac12\right)^T
\]
所以加入新变量后的表为
\[
\begin{array}{c|rrrrrr|r}
B & x_1&x_2&x_3&s_1&s_2&s_3&b\\
\hline
  & 0&0&-\frac32&0&\frac12&1&250\\
s_1&0&0&\frac12&1&\frac12&-1&50\\
x_2&0&1&\frac12&0&\frac12&0&100\\
x_1&1&0&\frac12&0&-\frac12&1&50
\end{array}
\]

因为 $x_3$ 的检验数是 $-\frac32<0$，故 $x_3$ 入基。
比值检验：
\[
\frac{50}{1/2}=100,\quad \frac{100}{1/2}=200,\quad \frac{50}{1/2}=100
\]
取 $x_1$ 出基。

第一次迭代后：
\[
\begin{array}{c|rrrrrr|r}
B & x_1&x_2&x_3&s_1&s_2&s_3&b\\
\hline
  & 3&0&0&0&-1&4&400\\
s_1&-1&0&0&1&1&-2&0\\
x_2&-1&1&0&0&1&-1&50\\
x_3&2&0&1&0&-1&2&100
\end{array}
\]

此时 $s_2$ 的检验数为 $-1<0$，故 $s_2$ 入基。
比值检验：
\[
\frac{0}{1}=0,\quad \frac{50}{1}=50
\]
故 $s_1$ 出基。

最终表：
\[
\begin{array}{c|rrrrrr|r}
B & x_1&x_2&x_3&s_1&s_2&s_3&b\\
\hline
  & 2&0&0&3&0&2&400\\
s_2&-1&0&0&1&1&-2&0\\
x_2&0&1&0&-1&0&1&50\\
x_3&1&0&1&1&0&0&100
\end{array}
\]

所有检验数都非负，故达到最优。

取非基变量
\[
x_1=s_1=s_3=0
\]
得
\[
x_2=50,\quad x_3=100
\]
所以
\[
w^*=-400
\]
原最大化问题最优值为
\[
z^*=400
\]

结论：
\[
(x_1,x_2,x_3)=(0,50,100),\qquad z^*=400
\]

\textbf{（2）}

新问题化为标准型：
\[
\min w=-x_1-2x_2
\]

s.t.
\[
\begin{aligned}
x_1+s_1&=100,\\
2x_2+s_2&=200,\\
x_1+x_2+s_3&=150,\\
4x_1+x_2+s_4&=250,
\end{aligned}
\qquad x_1,x_2,s_1,s_2,s_3,s_4\ge 0
\]






把新约束加到原最优表中。
由原最优表可得
\[
x_1=50+\frac12 s_2-s_3,\quad x_2=100-\frac12 s_2
\]
代入新约束
\[
4x_1+x_2+s_4=250
\]
得
\[
s_4+\frac32 s_2-4s_3=-50
\]
于是增广表为
\[
\begin{array}{c|rrrrrr|r}
B & x_1&x_2&s_1&s_2&s_3&s_4&b\\
\hline
  & 0&0&0&\frac12&1&0&250\\
s_1&0&0&1&\frac12&-1&0&50\\
x_2&0&1&0&\frac12&0&0&100\\
x_1&1&0&0&-\frac12&1&0&50\\
s_4&0&0&0&\frac32&-4&1&-50
\end{array}
\]

此时检验数都非负，但最后一行右端项为负，故改用对偶单纯形法。
由最后一行看，只有 $s_3$ 列为负，因此取 $s_3$ 入基，$s_4$ 出基。

一次对偶单纯形迭代后得到最终表：
\[
\begin{array}{c|rrrrrr|r}
B & x_1&x_2&s_1&s_2&s_3&s_4&b\\
\hline
  & 0&0&0&\frac78&0&\frac14&\frac{475}{2}\\
s_1&0&0&1&\frac18&0&-\frac14&\frac{125}{2}\\
x_2&0&1&0&\frac12&0&0&100\\
x_1&1&0&0&-\frac18&0&\frac14&\frac{75}{2}\\
s_3&0&0&0&-\frac38&1&-\frac14&\frac{25}{2}
\end{array}
\]

此时所有右端项都非负，且所有检验数都非负，所以达到最优。

取非基变量
\[
s_2=s_4=0
\]
得
\[
x_1=\frac{75}{2},\quad x_2=100
\]
故
\[
w^*=-\frac{475}{2}
\]
原最大化问题最优值为
\[
z^*=\frac{475}{2}
\]

结论：
\[
(x_1,x_2)=\left(\frac{75}{2},100\right),\qquad z^*=\frac{475}{2}
\]




\end{document}